\section{The \texttt{Lens\_parab} McXtrace Component}
X-ray compound refractive lens (CRL) with a profile of the parabola

\subsection*{Identification}
\begin{itemize}
  \item \textbf{Author:} Jana Baltser and Erik Knudsen
  \item \textbf{Origin:} NBI
  \item \textbf{Date:} August 2010, modified July 2011
\end{itemize}

\subsection*{Description}
A simple X-ray compound refractive lens (CRL) with a profile of the parabola in rotation simulates the photons' movement on passing through it. The CRL focuses in 2D

Example: Lens\_parab(material\_datafile = "Be.txt", r=200e-6, r\_ap=0.5e-3, d=50e-6, N=16)

\subsection*{Input parameters}
Parameters in \textbf{boldface} are required; the others are optional.

\begin{longtable}{p{0.22\textwidth}p{0.12\textwidth}p{0.46\textwidth}p{0.14\textwidth}}
\toprule
\textbf{Name} & \textbf{Unit} & \textbf{Description} & \textbf{Default} \\
\midrule
\endhead
material\_datafile & str & Datafile containing f1 constants & "Be.txt" \\
r & m & Radius of curvature (circular approximation at the tip of the profile). & 0.5e-3 \\
r\_ap & m & Radius of circular aperture, which also defines the depth of the lens profile. & 1.4e-3 \\
d & m & Distance between two surfaces of the lens along the propagation axis. & .1e-3 \\
N & m & Number of single lenses in a stack. & 1 \\
rough\_z & rad & RMS value of random slope error along z. & 0 \\
rough\_xy & rad & RMS value of random slope error along x and y. & 0 \\
\bottomrule
\end{longtable}

\subsection*{Links}
\begin{itemize}
  \item Component source code found in file \texttt{Lens\_parab.comp}.
\end{itemize}
\IfFileExists{optics/Lens_parab_static.tex}{\input{optics/Lens_parab_static.tex}}{}